\documentclass[11pt]{article}
\usepackage[margin=1in]{geometry}
\usepackage{booktabs}
\usepackage{url}
\title{Can the estimator's operating domain be computed?}
\date{Prepared October 5, 2026}
\begin{document}
\maketitle

\section{Question}
The user chose to enlarge the estimator's operating domain rather than limit
the controller, and asked whether that domain should be computed rather than
declared. This report classifies every domain premise of the NRMM estimator
(\path{config/nrmmTrackingConfig.m}, overridden by
\path{config/estimatorControllerIntegrationConfig.m}). For each premise it
records where the premise enters the design and the runtime, and whether it can
be derived from the controller's vehicle, tire and state-limit configuration.
The estimator is unchanged; the scripts are \path{scripts/studyEstimatorDomain.m}
and \path{scripts/replayEstimatorValidity.m}.

\section{Classification}
\begin{center}
\begin{tabular}{p{0.24\linewidth}p{0.16\linewidth}p{0.5\linewidth}}
\toprule
premise (declared) & role & derivation\\
\midrule
ego yaw-rate maximum (0.75~rad/s) & runtime check, caps & Not load-bearing: the velocity certificate $A_v^\top+A_v=-2\lambda I$ holds for every yaw rate. The controller's enforced limit $|r|\le5$~rad/s can be used.\\
ego speed interval (5--25~m/s) & gain, coefficients, kinematic check & From the controller's limits, 1--18~m/s.\\
single-track mismatch $|r-v_y/l_r|$ (0.25~rad/s) & velocity disturbance (dominant) & Equals $v_x|\tan\alpha_r|/l_r$. It is bounded only while the rear tire is below its Fiala saturation slip $\alpha_{\rm sat}=\arctan(3\mu F_z/C)=0.195$~rad, giving $\le2.15$~rad/s at 18~m/s.\\
ego sideslip inversion domain (0.12~rad) & kinematic inversion branch & Bounded by rear slip and yaw rate; no premise-free bound.\\
ego acceleration norm ($\infty$) & between-sample velocity bound & $\mu g$ plus road load: 8.52~m/s$^2$.\\
ego yaw acceleration (set by the driver) & orientation transport & Axle friction forces over $I_z$: 12.26~rad/s$^2$, already derived this way.\\
target speed, acceleration, sideslip & target model & A contract about other traffic; it cannot be derived from our vehicle. It can be taken from the scenario set.\\
relative position maximum (50~m) & gain, ranges & Radar range (sensor specification).\\
\bottomrule
\end{tabular}
\end{center}

\section{Cost of the computed values}
Design ultimate bounds from \texttt{synthesizeNrmmObserverGains}
(\path{ESTIMATOR_DOMAIN_20261005/ultimate-bounds.csv}):
\begin{center}
\begin{tabular}{lrrrr}
\toprule
domain & $\|e_v\|$ (m/s) & rel.\ position (m) & target vel.\ (m/s) & target acc.\ (m/s$^2$)\\
\midrule
declared & 0.488 & 1.480 & 8.06 & 23.3\\
yaw rate, speed from controller & 0.488 & 1.480 & 8.06 & 23.3\\
+ mismatch from rear saturation & 3.651 & 6.738 & 37.5 & 112.9\\
+ sideslip = rear saturation slip & 3.695 & 6.810 & 38.0 & 114.1\\
\bottomrule
\end{tabular}
\end{center}
Enlarging the yaw-rate and speed domains costs nothing. Computing the
single-track mismatch from the tire makes every guaranteed bound five to seven
times larger, which makes them useless for a 0.2-m margin.

\section{Validity with the free enlargement}
The estimator was replayed on the recorded noisy runs (aec20bc) with yaw-rate
maximum 5~rad/s and speed minimum 1~m/s
(\path{ESTIMATOR_DOMAIN_20261005/replay-enlarged-yaw-speed.csv}). At 15~m/s
the ego bound stays valid in all seven runs, against four before. At 8~m/s
it is still invalidated at 0.40--1.50~s, now by the kinematic feasibility
check. With the declared sideslip domain and mismatch, the measured lateral
velocity $l_r r$ must lie within $v\sin(0.12)+l_r(0.25+n_g)$. That is about
0.76~rad/s at 7~m/s, and the 8-m/s avoidances reach 0.81~rad/s. The vehicle's
real sideslip and mismatch exceed the declared premises there.

\section{Conclusion}
The premises that matter, mismatch and sideslip, bound body-point lateral
speed and sideslip (see the addendum: the estimator has no wheel model). The
plant alone does not bound them usefully; the controller can enforce them. The yaw-rate, speed, acceleration and
yaw-acceleration premises can be computed from the plant and the controller's
limits without loss. The target premises are a contract about other vehicles.

\section{Addendum: the estimator has no wheel model; premise audit}
The user asked whether the estimator models front and rear wheels. It does
not. The ego relation is purely kinematic (\path{estimator/OBSERVER_ISS_THEORY.md},
Sec.~3):
\[
\omega_E=v_y/l_E+d_{\rm st},\quad|d_{\rm st}|\le\delta_{\rm st},\qquad
v_x\ge\|v\|\cos b .
\]
$l_E$ (named \texttt{rearAxleDistance}) is a geometric length of the observer.
$v_y-l_E\omega_E$ is the lateral velocity of the body point $l_E$ behind the
centre of mass, so the premises say that this point's lateral speed is at
most $l_E\delta_{\rm st}$ (0.41~m/s with $l_E=1.65$~m) and that the sideslip is
at most $b=0.12$~rad. Tires enter only through the true vehicle. With
$l_E$ equal to the real rear-axle distance, that point is the rear axle, and
its lateral speed is $v_x\tan\alpha_r$. The earlier statement that the
estimator's premises concern the rear tire was an interpretation through the
plant, not a property of the estimator. $l_E$ is a free design length.

Recorded true motion (aec20bc): the point's lateral speed reached 0.42--0.58~m/s
at 8~m/s and 0.86--4.25~m/s at 15~m/s, and the sideslip reached 0.167~rad and
0.24~rad. The body point with the least lateral speed lies about 1.0~m behind
the centre of mass at 8~m/s (at most 0.10--0.24~m/s) and near the centre of
mass at 15~m/s.

The adapter's truth audit was counted over the replay with the yaw-rate and
speed domains enlarged (\path{ESTIMATOR_DOMAIN_20261005/audit-enlarged-yaw-speed.csv}).
The premises were violated in 6--93 frames per run. At 15~m/s, frames with an
available ego bound that did not contain the truth were: headOn 12 of 19,
acceleratingHeadOn 8 of 15, brakingLead 2 of 8, crossing 24 of 31,
turningCrossing 10 of 17, curvedHeadOn 24 of 31, curvedCrossing 12 of 20.
Target bounds failed in similar numbers. The runtime consistency check tests
only the measured lateral velocity against the forward cone, so these
violations pass undetected. At 8~m/s the bounds were available only in the
first 8--30 frames and always contained the truth.

The published bounds are therefore unsound in the 15-m/s avoidances, which
also affects every earlier noisy result that used them. Both premises are
affine or nearly affine in the controller's state [$v_y$, $r$, $v_x$], so the
controller could enforce them. Alternatively $l_E$ and $\delta_{\rm st}$ could
be chosen from the motion the controller is allowed to produce.
\end{document}
